\section{赛题理解与总体目标}
\label{sec:intro}

\subsection{赛题一任务拆解}

赛题一“欺诈交易智能识别”面向金融交易全链路风险防控场景，要求参赛团队自建或模拟多源异构交易环境，并在交易发生时完成实时风险评估、欺诈判定和风险等级输出。其目标不是单纯训练一个二分类模型，而是形成覆盖数据仿真、实时计算、模型识别、风险解释和持续迭代的完整反诈系统。

结合规则文件，赛题一的核心交付可以拆解为以下六类能力：

\begin{enumerate}
  \item \textbf{仿真环境：}模拟用户、账户、商户、设备、IP、渠道、支付方式等对象，生成正常交易与多类型欺诈交易，体现跨时段和跨渠道行为；
  \item \textbf{实时检测：}将交易组织为流式输入，建立数据接入、画像补齐、窗口特征计算、模型推断、风险输出的在线链路；
  \item \textbf{多维特征：}融合用户行为、设备指纹、IP环境、商户画像、短时交易窗口、交易网络和图挖掘特征；
  \item \textbf{模型迭代：}支持多模型训练、阈值校准、人工反馈样本回流、轻量重训和在线模型热切换；
  \item \textbf{团伙识别：}基于共享设备、共享IP、共同收款方和高风险商户挖掘疑似欺诈团伙，并输出可解释证据；
  \item \textbf{工程落地：}提供可部署架构、性能测试方案、监控看板、测试用例和安全合规设计。
\end{enumerate}

\subsection{作品定位}

本作品命名为FP-FraudSim流式风控系统，定位为面向第三方支付和互联网金融交易的实时欺诈识别原型。系统以“仿真交易流 + 流式风险计算 + 模型服务 + 图谱溯源 + 人工反馈”为主线，覆盖赛题一完成度评分中的四个明确子项：仿真环境与实时检测能力、流式数据处理框架、自适应学习与模型迭代能力、多维度特征构建。

与只提交离线Notebook或静态分类结果的方案相比，本系统强调三个工程闭环。第一，数据闭环：通过ETL、七类欺诈注入脚本与可调参数仿真器生成可复现实验流，再由Producer按速率回放至Kafka。第二，识别闭环：Flink计算窗口特征并调用模型服务，风险结果实时进入Kafka主题和Dashboard。第三，迭代闭环：人工复核标签同时写入反馈主题和样本池，训练脚本生成候选模型，经指标对比后受控发布，并保留快速回滚版本。

\subsection{问题驱动的技术路线}

金融交易反诈具有\textbf{风险发生快、欺诈样本少、行为关系复杂、业务误报成本高、欺诈手法持续变化}等特点。因此，技术路线不能只追求离线分类精度，而需要同时解决实时性、识别效果、可解释性和持续迭代问题。

\begin{enumerate}
  \item \textbf{针对“事后识别来不及”的问题，选择Kafka + Flink流式架构。}Kafka将交易接入与风险计算解耦，可缓冲流量峰值并支持结果回放；Flink面向有状态流计算，能够按事件时间维护用户、设备和商户窗口，并通过watermark处理乱序交易。这比定时批处理更适合在支付发生时完成风险判定。
  \item \textbf{针对“单笔交易看似正常”的问题，构建画像、窗口与图关联多维特征。}用户历史画像用于判断行为偏离，短时窗口用于识别爆发式异常，设备/IP和图关联用于发现多账户协同关系，使系统从单笔静态判断升级为跨时段、跨实体研判。
  \item \textbf{针对“结构化数据、类别特征多且样本不均衡”的问题，选择梯度提升模型作为主模型。}LightGBM、XGBoost和CatBoost能够有效处理非线性特征交互、缺失值及类别变量，在中等规模表格数据上训练和推理成本较低；系统以PR-AUC、Top 1\%召回和误报率评价模型，更符合低欺诈率业务场景。
  \item \textbf{针对“团伙欺诈难以由单笔模型发现”的问题，引入可解释图挖掘和GraphSAGE旁路实验。}共享资源与连通分量方法负责稳定产出证据链；GraphSAGE验证图嵌入价值，但不直接影响线上主模型，兼顾创新性与稳定性。
  \item \textbf{针对“黑盒结果难以直接处置”的问题，输出理由码并提供阈值沙盒。}模型分数被映射为放行、复核和拒绝决策，评审或业务人员可调整两级阈值，实时观察精度、召回、误拦和审核工作量变化。
  \item \textbf{针对“欺诈模式持续漂移”的问题，设计受控模型生命周期。}人工反馈进入样本池后生成candidate候选模型，经评估后提升为latest，旧模型自动备份为rollback，避免反馈驱动重训直接覆盖线上模型。
\end{enumerate}

\subsection{评分准则对应关系}

表~\ref{tab:score_mapping}按照赛题一评分细则拆分系统能力，并列出评审时可以直接核验的交付证据。后文按照该映射展开技术实现、实验验证和安全设计。

\begin{table}[H]
  \centering
  \caption{赛题评分点与系统实现对应关系}
  \label{tab:score_mapping}
  \renewcommand\arraystretch{1.3}
  \setlength{\tabcolsep}{3pt}
  {\songti\zihao{5}
  \begin{tabularx}{0.98\textwidth}{>{\raggedright\arraybackslash}p{1.45cm}>{\centering\arraybackslash}p{0.7cm}>{\raggedright\arraybackslash}X>{\raggedright\arraybackslash}X}
    \toprule
    评分项 & 分值 & 核心实现 & 可核验交付证据 \\
    \midrule
    仿真与实时检测 & 5 & 七类场景注入、可调参数仿真、交易流回放与风险分级 & 仿真配置、Producer、实时结果与告警界面 \\
    流式处理框架 & 5 & Kafka接入、Flink事件时间与窗口计算、迟到事件分流 & Docker服务、Flink UI、风险/告警/迟到主题 \\
    自适应学习与迭代 & 5 & 反馈样本池、candidate/latest/rollback、阈值沙盒 & 反馈接口、候选指标、发布/回滚与重载接口 \\
    多维特征构建 & 5 & 交易、画像、设备/IP、窗口、商户和图关联特征 & 特征配置、窗口代码、Redis画像、图特征产物 \\
    创新性 & 20 & 微批评分、图证据、GraphSAGE旁路、阈值沙盒 & 图实验指标、证据图、阈值交互、一键演示 \\
    实用性 & 20 & 风险处置、接口接入、边云部署与性能验证 & API输出、Dashboard、Compose/K8s清单、压测报告 \\
    成熟度 & 20 & 模块化、多模型适配、测试覆盖、受控模型生命周期 & 单元测试、版本产物、发布回滚接口、技术文档 \\
    安全性 & 20 & API Key、双通道审计、Secret与网络隔离 & audit日志/主题、K8s Secret、NetworkPolicy、安全方案 \\
    \midrule
    技术作品项合计 & 100 & 覆盖完成度、创新性、实用性、成熟度和安全性 & 由系统演示、代码工程、测试结果和本文档共同证明 \\
    \bottomrule
  \end{tabularx}
  }
\end{table}

\noindent\textbf{说明：}团队表现为额外10分加分项，主要通过答辩逻辑、现场协作和问题应答体现，不计入表~\ref{tab:score_mapping}所列技术作品项100分。

\subsection{文档结构}

本文档第二节说明数据仿真、欺诈场景和特征体系；第三节介绍系统架构、实时处理链路、模型训练、图挖掘和可视化交互；第四节给出实验验证、性能测试方案与复现流程；第五节总结创新点、应用价值、安全合规和后续优化方向。
